% ------------------------------------------------------------------
%  Conclusion generale (questions V.1 a V.4 des consignes)
% ------------------------------------------------------------------
\section*{Conclusion générale}
\phantomsection
\addcontentsline{toc}{section}{Conclusion générale}
\markboth{Conclusion générale}{}

\paragraph{Utilité de la salle blanche.}
En microfabrication, une particule posée sur le wafer est recopiée dans le motif. Une particule d'un micromètre peut relier deux pistes, couper une piste de \qty{2}{\micro\metre} ou percer un film mince, et le rendement chute quand la densité de particules augmente. La salle blanche filtre l'air et l'envoie en flux descendant pour évacuer les particules. Elle règle aussi la température et l'humidité, dont dépendent l'épaisseur des résines et leur photoréaction, et utilise un éclairage jaune qui n'expose pas les résines \cite{intro,cours303}. L'opérateur étant la principale source de particules, la tenue complète est obligatoire.

\paragraph{Classe d'une salle blanche.}
La norme ISO 14644-1 classe une salle selon la concentration maximale de particules en suspension dans l'air \cite{cours303}. Une salle ISO $N$ admet au plus $10^N$ particules de taille $\geq \qty{0,1}{\micro\metre}$ par mètre cube. Une salle ISO 5 admet ainsi \num{100000} particules $\geq \qty{0,1}{\micro\metre}$ et \num{3520} particules $\geq \qty{0,5}{\micro\metre}$ par \unit{\metre\cubed}, l'équivalent de l'ancienne classe américaine 100. Plus le chiffre est petit, plus la salle est propre : au CMi, le niveau $-1$ est classé ISO 5 et le niveau +1 ISO 6 ou ISO 7 selon les zones.

\paragraph{Relation entre les sujets.}
Les six sujets forment la chaîne de fabrication d'une même résistance. Le sujet~1 calibre la lithographie de l'AZ1512 (épaisseur, dose), résine réutilisée au sujet~4 ; les sujets~2 et 5 emploient d'autres résines, réglées d'après les fiches techniques et les recettes du CMi. Les sujets~2, puis 4 et 5, sont les deux façons de structurer l'Al : le lift-off et la gravure humide. Le sujet~3 mesure la géométrie obtenue : undercut et largeur optique des pistes. Le sujet~6 mesure la résistance, en tire $R_\square$ et la largeur effective $W$ des pistes par TLM, puis $\rho$ avec l'épaisseur $T$ mesurée au sujet~5 (équation~\ref{eq:resistance}).

\paragraph{Nos résultats et le TP.}
Les résultats principaux sont cohérents d'un sujet à l'autre : l'épaisseur d'Al mesurée au Dektak (\qty{554}{\nano\metre}) donne au sujet~6 une résistivité à \qty{4}{\percent} de l'Al massif, et la comparaison des deux procédés est nette, le lift-off respectant la largeur du masque alors que la gravure humide fait perdre \qty{2,5}{\micro\metre} à chaque piste. Le TP nous a surtout appris l'importance de la méthode. Au sujet~1, exposer les quatre wafers en parallèle nous a privés de la courbe de contraste des films minces ; il aurait fallu mesurer un wafer avant d'exposer le suivant. Certaines données manquent : la durée de gravure au CMi n'a pas été relevée, et le fichier de nos mesures du sujet~6, prises pendant la séance, a été perdu ; nous avons donc utilisé celles du groupe~10. Cela limite certaines conclusions : noter systématiquement temps et paramètres sur place, et sauvegarder les fichiers à plusieurs endroits, est la première leçon que nous en tirons. Enfin, la séance au CMi a montré ce qu'apportent la salle blanche et l'automatisation en propreté et en reproductibilité, par rapport au travail manuel en DLL-MED.
